\chapter{Software and Reproducibility}
\label{app:software}

The methods in this book are implemented, open source, and testable, and the exercises
marked \emph{Implementation} assume the reader will write code. This appendix says what
exists, what each package is for, and how to check an implementation against something
other than itself.

\section{Reference implementations}

\paragraph{Dual active-set QP.} The method of Chapter~\ref{ch:walking}, together with
the guessing fast path of Chapter~\ref{ch:guessing} and the factor reuse of
Chapter~\ref{ch:warmstart}, is implemented in the \texttt{cvx-quadprog} package
\cite{cvxquadprog}. Its signature carries the two conventions this book inherited: the
linear term is subtracted and constraints are held column-wise.

\paragraph{Non-negative conjugate gradients.} The loop of Chapter~\ref{ch:guessing} with
the matrix-free Krylov inner solver of Chapter~\ref{ch:krylov}, including the
equality-augmented variant, the regularising split, and the Nyström preconditioners of
Chapter~\ref{ch:conditioning}, is released as \texttt{nncg}
\cite{schmelzer2026nncg}. Its continuous-integration suite is the paper's synthetic
study, which is a good model to copy: \emph{the test suite is the experiment}.

\paragraph{The Critical Line Algorithm and the path tracer.} The homotopy of
Chapter~\ref{ch:parametric}, the operator contract of Chapter~\ref{ch:operator} with its
four backends, and the LASSO client of Chapter~\ref{ch:homotopy} are implemented in
\texttt{cvxcla} \cite{cvxcla}. It factors into a problem-agnostic path tracer of which
the frontier tracer and the LASSO solver are two clients---which is the software form of
Theorem~\ref{thm:identity}.

\paragraph{General-purpose comparators.} Clarabel \cite{clarabel2024} (interior point)
and OSQP \cite{osqp} (operator splitting) are the solvers Chapter~\ref{ch:measurement}
benchmarks against, and both have native APIs worth calling directly. CVXPY
\cite{diamond2016} is the modelling layer whose overhead that chapter warns about
attributing to an algorithm.

\section{How to validate an implementation}

Four levels, in increasing order of how much they tell you. Every exercise in this book
marked \emph{Implementation} can be checked at one of them.

\paragraph{1. The certificate.} Every returned point must satisfy the KKT
test~\eqref{eq:certificate} \emph{against the original data}, not against the
construction that produced it. This is cheap, it is self-contained, and it catches most
errors. Make it an assertion in your test suite, not a diagnostic you run when
suspicious.

\paragraph{2. An independent solver.} Solve the same instance with a general-purpose QP
solver and compare minimisers. Agreement to solver precision---typically $10^{-8}$ or
better on well-conditioned data---is the expectation. For a parametric method, compare
at each of a grid of parameter values: the traced segment evaluated at a grid point
should reproduce the independently solved QP there.

\paragraph{3. A closed form the code does not use.} Chapter~\ref{ch:walking} gives the
model: the iteration computes $z$ and $r$ from carried factors, and
Proposition~\ref{prop:invariance} gives closed forms in $(G,A,n_*)$ alone. Computing
both and comparing to machine precision tests the factorisation machinery against the
mathematics rather than against a second implementation of itself. Similarly,
Proposition~\ref{prop:increment} lets a running objective be checked against a fresh
evaluation at every step.

\paragraph{4. A different field's reference.} The most convincing check available for
Part~III, and one that only the identity makes possible: trace a LASSO path with a
frontier tracer and compare against an established LARS implementation, breakpoint by
breakpoint. Agreement at the $10^{-15}$ level on a standardised design, including on
high-dimensional problems, is achievable and is a genuinely independent test---different
authors, different field, different code lineage.

\section{Two hazards worth building tests for}

\paragraph{The hazard that never fires on your test family.} Chapter~\ref{ch:measurement}
made the general point; here is the concrete test. Duplicate a column of the constraint
matrix and re-run your guessing fast path. If your solver has no rank test and no
certificate fallback, it will now return plausible wrong answers, and it will do so on
the sort of model any programmatic assembly can produce. Add the test.

\paragraph{The degeneracy that only appears at scale.} The feasible-set projection of
Chapter~\ref{ch:numerics}, Section~6, addresses a failure that appears only after
hundreds of steps on a near-rank-deficient operator. A test suite that runs only small
well-conditioned instances will never see it. Build a controlled rank-deficiency sweep
and assert both branches: that the trace completes and matches a reference while the free
block stays full rank, and that the solver \emph{declines} once it does not.

\section{On reproducing numbers}

Every empirical figure quoted in this book comes from the companion papers, and each is
reproducible from their experiment code---the S\&P~500 and FTSE~100 studies, the
synthetic scaling studies, the constraint-family sweeps behind the $93\%$ figure of
Chapter~\ref{ch:guessing}, and the out-of-sample study of Chapter~\ref{ch:measurement}.

Two habits from those papers are worth copying regardless of subject. First, the figures
quoted in the text are generated by the same run that produces the tables, so prose and
table cannot drift apart. Second, empirical claims made anywhere in a paper are measured
in its own experiments rather than quoted from elsewhere. Both are cheap to adopt and
they eliminate an entire category of error that peer review does not catch.
